\section{Implementation}\label{sec:implementation}
%
Our Java implementation of the boundedness analysis is
freely available, with instructions online\footnote{\url{https://github.com/spoto/boundedness-analysis/blob/main/README.md}} on how to run and verify
the experiments in this paper. It issues a warning in \<Ponzi1> and \<Lottery1>
and no warning in \<Ponzi2> and \<Lottery2>, as expected.

We do not know of any other boundedness analysis that we can compare to ours.
There are, however, tools that compute the gas consumption of the methods
of a smart contract, as a formula on the size of some variables.
The problem is different but there is an intersection:
if a method has a constant complexity, then all iterations inside it
must be bounded. In particular, we have run the GASTAP tool
on our examples. Originally, it was the gas cost inference engine of
GASOL~\cite{AlbertCGRR20}, but it is now available as a stand-alone
tool accessible through a web interface\footnote{\url{https://costa.fdi.ucm.es/gastap/}}. GASTAP requires smart contracts written in Solidity $0.4.24$,
hence we have translated our examples\footnote{\url{https://github.com/spoto/boundedness-analysis/tree/main/contracts}}.
The only difficulty here is
that Solidity has no \<removeFirst()> method, hence the translation
of \<Ponzi2> uses an extra variable \<start> that refers to the valid portion of the
list \<investors>.
GASTAP reports a gas consumption for method \<invest()> of \<Ponzi1> of
$55160+6558*\|nat|(\<investors>)$\ie linear in the number of investors,
which correctly cannot exclude an unboundedness issue in \<Ponzi1>.
The gas consumption of method \<invest()> of \<Ponzi2> is reported to be
$717323+28488*\|nat|(\<investors>-\<start>-100)$, which correctly corresponds
to the cost of that method but which is not constant,
hence nothing can be concluded from that about the boundedness of the loop
at line~\ref{ponzi2:iteration} of \<Ponzi2> (Fig.~\ref{fig:ponzi2}).
The gas comsumption for method \<drawWinner()> of \<Lottery1> and \<Lottery2>
is the same: $11558+94*c(\|maximize\_failed|)+50*\|nat|(-\|unknown|+32)+6*(\|nat|(\|unknown|)/32)+\|failed|(\|no\_rf|)$.
Beyond the interpretation of that expression,
what is relevant is that it is not a constant, hence
nothing can be concluded about the boundedness of the loop
at line~\ref{lottery1:iteration} in \<Lottery1> (Fig.~\ref{fig:lottery1})
and in \<Lottery2>, while
we know that that loop is bounded in \<Lottery2> and unbounded in \<Lottery1>.
Again, this just shows that inference of gas costs and boundedness analysis
have an intersection, but they are different problems that expect different
solutions.
